\subsection{Sägezahn-Schwingung}
Gleichung \ref{eq:Saegezahn} zeigt die Fourier-Zerlegung einer Sägezahnschwingung bekannt aus der Vorlesung und Übung.
\begin{equation}
    \label{eq:Saegezahn}
    f = - \frac{2}{\pi}\sum_{k\geq1}\frac{(-1)^k}{k}\sin(kx)
\end{equation}
Aus der Gleichung folgt, das eine Sägezahnschwingung aus einer Grundfrequenz und deren ganzahligen Vielfachen zusammengesetzt wird. Wird auf eine solche Schwingung nun ein sogenannter Tiefpassfilter angewendet, dessen Filterfrequenz sich der Grundfrequenz der Sägezahnschwingung annähert $(f_{filter} \to x)$, so ensteht wiederum eine reine harmonische Schwingung. Aufgrund des unterschiedlichen Klangbilds einer Sägezahnschwingung im Vergleich zu einer rein harmonischen Schwingung kann dieser Effekt ausgenutzt werden, um den Klang eines Synthesizers anzupassen.\\
Die durch einen Synthesizer erzeugte Sägezahnschwingung wird hierbei durch einen Tiefpassfilter geleitet, um so höher die eingestellte Filterfrequenz ist, desto mehr gleicht die resultierende Schwingung dem Original, der Ton klingt also, typisch für Sägezahnschwingungen, rau, $f_{filter} \to \infty$ wäre hierbei die perfekte Sägezahnschwingung. Je mehr die Filterfrequenz nun der Grundfrequenz des erzeugten Tons angenähert wird, desto weicher klingt der resultierende Ton.